\section{Tổng quan dự án}
\label{sec:overview}

\subsection{Giới thiệu bài toán}

Dự án này được lấy cảm hứng từ cuộc thi Kaggle \textit{``Playground Series --- Season 4, Episode 11: Exploring Mental Health Data''}~\cite{kaggle2024}. Mục tiêu chính là xây dựng mô hình học máy để dự đoán liệu một cá nhân có nguy cơ bị trầm cảm (\textit{Depression}) hay không, dựa trên dữ liệu khảo sát sức khỏe tâm thần.

\textbf{Phát biểu bài toán:} Cho trước một tập dữ liệu gồm các đặc trưng nhân khẩu học, học vấn, nghề nghiệp và các yếu tố tâm lý --- xã hội, bài toán yêu cầu phân loại nhị phân (\textit{Binary Classification}) để dự đoán biến mục tiêu \texttt{Depression} nhận giá trị 0 (không trầm cảm) hoặc 1 (trầm cảm).

\subsection{Mô tả dữ liệu}

Bộ dữ liệu được cung cấp bởi Kaggle với các thông tin chính:

\begin{itemize}[noitemsep]
    \item \textbf{Tập huấn luyện (train)}: 140.700 mẫu với 20 đặc trưng và 1 biến mục tiêu
    \item \textbf{Tập kiểm tra (test)}: 93.800 mẫu (không có nhãn)
    \item \textbf{Tính chất}: Dữ liệu tổng hợp (\textit{synthetic}), được tạo từ mô hình deep learning huấn luyện trên dữ liệu khảo sát thực tế
    \item \textbf{Metric đánh giá}: Accuracy (theo yêu cầu cuộc thi)
\end{itemize}

\subsection{Pipeline tổng thể}

Hình~\ref{fig:pipeline} minh họa toàn bộ quy trình xử lý dữ liệu và xây dựng mô hình của dự án, bao gồm 6 giai đoạn chính từ dữ liệu thô đến kết quả dự đoán cuối cùng.

\begin{figure}[H]
\centering
\begin{tikzpicture}[
    node distance=1.2cm,
    block/.style={rectangle, rounded corners, draw=black!70, fill=blue!8,
                  text width=3.2cm, minimum height=1cm, align=center, font=\small\bfseries},
    arrow/.style={-{Stealth[length=3mm]}, thick, draw=black!60}
]
\node[block] (raw) {1. Dữ liệu thô\\(Raw Data)};
\node[block, right=of raw] (eda) {2. Khám phá\\dữ liệu (EDA)};
\node[block, right=of eda] (prep) {3. Tiền xử lý\\(Preprocessing)};
\node[block, below=of prep] (fe) {4. Kỹ thuật\\đặc trưng (FE)};
\node[block, left=of fe] (model) {5. Xây dựng\\mô hình};
\node[block, left=of model] (eval) {6. Đánh giá\\(Evaluation)};

\draw[arrow] (raw) -- (eda);
\draw[arrow] (eda) -- (prep);
\draw[arrow] (prep) -- (fe);
\draw[arrow] (fe) -- (model);
\draw[arrow] (model) -- (eval);
\end{tikzpicture}
\caption{Pipeline tổng thể của dự án dự đoán rủi ro trầm cảm}
\label{fig:pipeline}
\end{figure}
